% apendix
\section{Historischer Kontext}

\subsection{Deutscher Orden}

Der \textit{Ordo Teutonicus} -- der Deutsche Orden -- wurde während der Belagerung von Akkon (1189--1191) als Krankenpflege-Bruderschaft deutscher Pilger gegründet und im März 1198 auf Betreiben des Bischofs Wolfger von Erla zum Ritterorden erhoben. Papst Innozenz~III. bestätigte diese Erhebung am 19.~Februar 1199. Erster Hochmeister war Heinrich Walpot von Bassenheim. Vorbild für Ordensregel und Struktur waren Johanniter und Templer; wie diese legten die Mitglieder die drei monastischen Gelübde ab und traten damit als \textit{milites Christi} in den Dienst des Heiligen Landes. \cite{militzer_deutscher_orden,boockmann_deutscher_orden}

\subsubsection*{Die Goldbulle von Rimini und der Beginn der Baltikum-Feldzüge (1226)}

Im Jahr 1226 wandte sich der polnische Herzog Konrad~I. von Masowien an den Deutschen Orden um Hilfe in seinem Krieg gegen die heidnischen \textit{Prußen} östlich der Weichsel. Hochmeister Hermann von Salza -- einer der einflussreichsten Diplomaten seiner Zeit -- ließ sich die Bedingungen vorab kaiserlich verbrieven: Kaiser Friedrich~II. stellte mit der \textbf{Goldenen Bulle von Rimini} die rechtliche Grundlage dafür, dass das nach der Missionierung und Unterwerfung der Prußen eroberte Land dauerhaft an den Orden fallen sollte. Papst Gregor~IX. bestätigte dies kurz darauf mit der Bulle von Rieti. Der Orden sicherte sich damit eine einzigartige Stellung: Er war künftig nur dem Papst, keinem weltlichen Lehnsherrn, unterstellt. \cite{militzer_deutscher_orden,boockmann_deutscher_orden}

1230 übergab Konrad~I. dem Orden im Vertrag von Kruschwitz das Kulmerland auf ewige Zeiten. Im Jahr 1231 überquerte Landmeister Hermann von Balk mit sieben Ordensrittern und etwa 700 Mann die Weichsel und errichtete im Kulmerland die erste Ordensburg \textit{Thorn}. Von dort aus begann die schrittweise Unterwerfung des Prußenlandes, begleitet von Christianisierung und systematischer Besiedlung. Die Teilnehmer dieser Feldzüge erhielten von Papst Gregor~IX. denselben Ablassbrief wie Kreuzfahrer ins Heilige Land. \cite{militzer_deutscher_orden}

\subsubsection*{Eingliederung des Schwertbrüderordens (1236/1237)}

Der 1202 in Riga gegründete \textit{Schwertbrüderorden} -- erkennbar am weißen Mantel mit rotem Schwert-Kreuz-Zeichen -- operierte eigenständig in Livland. Im Jahr 1236 erlitt er in der \textbf{Schlacht von Schaulen (Saule)} eine vernichtende Niederlage gegen schamai\-tische Litauer und Semgaller, bei der der Großteil seiner Ritter fiel. Hochmeister Hermann von Salza handelte daraufhin persönlich in Viterbo mit der päpstlichen Kurie die Eingliederung des geschwächten Ordens aus. Die \textit{Union von Viterbo} (1237) gliederte die verbliebenen Schwertbrüder als \textit{Livländischen Orden} in den Deutschen Orden ein. Damit verfügte der Orden nun über zwei Kernterritorien: Preußen im Westen und das Meistertum Livland im Osten. \cite{benninghoven_schwertbrueder,selart_livland}

\subsubsection*{Die Niederlage auf dem Peipussee (1242)}

In den 1240er Jahren stießen die livländischen Ordensritter, unterstützt vom Bischof von Dorpat (Hermann~I. von Buxthoeven), von Estland aus in das Gebiet der \textit{Republik Nowgorod} vor. 1240 gelang ihnen die Besetzung der strategisch wichtigen Stadt Pskow. Der Nowgoroder Fürst \textbf{Alexander Jaroslawitsch} -- bald \textit{Newski} genannt nach seinem Sieg an der Newa im selben Jahr über ein schwedisches Heer -- wurde zum Gegner. Im März 1242 befreite er Pskow zurück und trug den Kampf ins Ordensland. \cite{selart_livland}

Am \textbf{5.~April 1242} kam es auf dem zugefrorenen \textit{Peipussee} zur entscheidenden Begegnung. Die Ordensritter griffen in ihrer bewährten Keilformation gegen das Zentrum des Rus'-Heeres an. Newski hatte seine Druschina -- die schwere Reitertruppe -- im Hinterhalt zurückgehalten. Als die Ritter das Fußvolk fast durchbrochen hatten und das abschüssige Ufer ihre Schlagkraft lähmte, schickte Newski seine Kavallerie um den rechten Flügel der Ordensritter und traf sie von hinten. Eingekesselt und ihrer Taktik beraubt, wurden viele Ritter auf dem Eise erschlagen; überliefert ist auch, dass das Eis unter dem Gewicht der Panzerreiter stellenweise nachgab. Russischen Quellen zufolge wurden rund 500 Kämpfer getötet und 50 adlige Ordensritter gefangen; westliche Historiker gehen von deutlich niedrigeren Zahlen aus. \cite{selart_livland,benninghoven_schwertbrueder}

Im Sommer 1242 schloss der Orden Frieden mit Nowgorod. Die \textit{Narwa} wurde als Grenzfluss festgelegt; der Orden verzichtete dauerhaft auf weitere Expansion in nowgorodisches Gebiet. Die Ostgrenze des Ordensstaates war damit für mehr als ein Jahrhundert definiert. Die historische Bedeutung der Niederlage wurde in russischen Chroniken bereits im Mittelalter zur Entscheidungsschlacht zwischen römisch-katholischer Kirche und russischer Orthodoxie überhöht. \cite{selart_livland}

\subsubsection*{Ordensverfassung und Alltagsleben}

Der Orden gliederte sich in drei Stände: \textit{Ritterbrüder} (die militärische Kraft, Laien), \textit{Priesterbrüder} (Liturgie, Kanzlei, Chronistik) und \textit{dienende Halbbrüder} (Haushalt, Wachen). Mehrere Kommenden bildeten eine \textit{Ballei}, geleitet von einem Landkomtur. An der Spitze stand der Hochmeister, unterstützt von den Großgebietigern (Großkomtur, Ordensmarschall, Großspittler, Ordenstrappier, Ordenstressler). Der Tagesablauf war streng geregelt: Gebet, Waffendienst, Verwaltung. Das Tragen eines weißen Mantels mit schwarzem Kreuz kennzeichnete die Ritterbrüder im Feld wie im Kloster. \cite{boockmann_deutscher_orden,militzer_deutscher_orden}

\subsection{Das Ende der Babenberger und Niederösterreich}\label{ende_der_babenberger}

Die \textit{Babenberger} herrschten seit 976 als Markgrafen, seit 1156 als Herzöge über Österreich. Unter \textbf{Leopold~VI. dem Glorreichen} (reg. 1198--1230) erlebte das Herzogtum eine kulturelle Blüte: Wien wurde zu einem Zentrum des Minnesangs, Klöster und Städte wurden gegründet, das Rechtswesen ausgebaut. Leopold~VI. nahm an Kreuzzügen teil -- darunter dem Fünften Kreuzzug nach Ägypten (1217--1221) -- und vermittelte in späteren Jahren als Diplomat zwischen Kaiser Friedrich~II. und dem Papsttum. \cite{lechner_babenberger,niederstaetter_oesterreich}

Sein Sohn \textbf{Friedrich~II. der Streitbare} (reg. 1230--1246) brach mit dieser Tradition diplomatischer Klugheit. Er konfiszierte Güter seiner Mutter und Schwester, führte Krieg gegen ungarische und böhmische Nachbarn und geriet in offenen Konflikt mit Kaiser Friedrich~II. -- was 1236 zur Verhängung der Reichsacht und seiner vorübergehenden Vertreibung aus dem Herzogtum führte. Die Exkommunikation des Kaisers durch den Papst eröffnete Friedrich kurzzeitig die Chance, die Erhebung Österreichs zum Königreich auszuhandeln; die Verhandlungen blieben jedoch ohne Ergebnis. \cite{lechner_babenberger,niederstaetter_oesterreich}

\subsubsection*{Die Nachfolgekrise (1246--1278)}

Am \textbf{15.~Juni 1246} fiel Friedrich~II. in der \textbf{Schlacht an der Leitha} gegen ungarische Truppen König Bélas~IV. Er starb ohne männliche Nachkommen -- das Haus Babenberg erlosch damit im Mannesstamm. Das \textit{Privilegium minus} von 1156 hatte zwar auch weibliche Erbfolge zugelassen, doch die Nachfolge blieb umstritten. \cite{lechner_babenberger,niederstaetter_oesterreich}

Friedrichs Nichte \textbf{Gertrud} (Tochter seines bereits 1228 verstorbenen Bruders Heinrich) heiratete noch 1246 Markgraf Vladislav von Mähren -- einen Sohn des böhmischen Königs Wenzel~I. --, doch Vladislav starb bereits am 3.~Januar 1247. In zweiter Ehe heiratete sie 1248 den badischen Markgrafen \textbf{Hermann~VI.}, der 1250 ebenfalls verstarb. Der dieser Ehe entstammende Sohn Friedrich führte bis zu seinem Tod 1268 den Titel \textit{Herzog von Österreich und Steier}, konnte aber die Herrschaft nie tatsächlich ergreifen. \cite{lechner_babenberger}

Der österreichische Adel wandte sich 1251 an den böhmischen König, dessen Sohn \textbf{Ottokar~II. Přemysl} Ende 1251 in Österreich einmarschierte und 1252 Friedrichs Schwester \textbf{Margarete} heiratete -- und damit formell Herzog von Österreich wurde. Ottokar~II. folgte 1253 seinem Vater auf dem Böhmen-Thron und herrschte in der Folge über ein mächtiges Gebilde von Böhmen bis Kärnten. Er unterlag jedoch \textbf{Rudolf~I. von Habsburg} in der \textbf{Schlacht auf dem Marchfeld} (26.~August 1278), wo er den Tod fand. Das Erbe der Babenberger fiel damit an das Haus Habsburg -- eine Herrschaft, die Österreich bis 1918 prägen sollte. \cite{lechner_babenberger,niederstaetter_oesterreich}

\subsubsection*{Burg Ottenstein im Waldviertel}\label{burg_ottenstein}

Im \textit{Waldviertel}, dem waldreichen Hochland nördlich der Donau, stand und steht die \textbf{Burg Ottenstein} -- eine Höhenburg auf einer Anhöhe über dem Kamptal (heute über dem Stausee Ottenstein). Einer der frühen Besitzer, \textbf{Hugo de Ottenstaine}, ist 1177 urkundlich erstmals belegt; die Burg selbst dürfte noch älter sein. Die Herren von Ottenstein waren verwandt mit den Herren von Rauheneck bei Baden. Die \textit{Burgkapelle}, deren Fresken aus der zweiten Hälfte des 12.~Jahrhunderts stammen und erst 1975 freigelegt wurden, zeugt von der frühen Bedeutung des Sitzes. \cite{noe_burgen_online_ottenstein,reichhalter_kuehtreiber,burgen_austria_ottenstein}

Im Hochmittelalter war das Waldviertel ein Grenzland, dessen Burgherren sowohl als Verwaltungsträger der Herzöge als auch als eigenständige Lokalmächte agierten. Die Burg Ottenstein blieb bis in die erste Hälfte des 15.~Jahrhunderts im Besitz der Ottensteiner und bildete damit über Jahrhunderte einen stabilen Herrschaftsmittelpunkt inmitten der bewaldeten Hochfläche. \cite{reichhalter_kuehtreiber,noe_burgen_online_ottenstein}

\subsubsection*{Die Herren von Ottenstein -- Stammbaum und Genealogie}\label{herren_von_ottenstein}

Die urkundliche Überlieferung zur Familie der \textbf{Herren von Ottenstein} ist fragmentarisch. Der folgende Stammbaum kombiniert gesicherte Belege mit historisch plausiblen Ergänzungen; letztere sind als \textit{[fiktional]} gekennzeichnet.

\medskip
\noindent\textbf{Quellen:} \cite{noe_burgen_online_ottenstein,reichhalter_kuehtreiber,schneider_ottenstein,burgen_austria_ottenstein}

\medskip
\noindent\textbf{Stammbaum (soweit belegbar):}

\begin{description}
  \item[\textbf{Otto de Staine} (\textasciitilde{}1156/71)] Vermutlicher Erbauer der Burg Ottenstein; neuere Forschungen (Kupfer) setzen die Errichtung der Anlage unter seiner Hand an. Nach Kupfer unmittelbar verwandt mit den \textit{Herren von Rauheneck} (Ministerialengeschlecht bei Baden); nach Marian hingegen ein landesfürstlicher Ministeriale aus der Sippe der Herren von Winkl. Der Name weist auf einen \textit{Otto} als Gründer hin. \cite{noe_burgen_online_ottenstein}

  \item[\textbf{Hugo de Ottenstaine} (\textasciitilde{}1177--\textasciitilde{}1192)] Erster namentlich belegter Burgherr. Zwischen 1177 und 1192 mehrfach urkundlich nachweisbar; Schenkungen an das \textit{Stift Zwettl} bezeugen Wohlstand und religiöses Engagement. Die Burg war freies Eigen der Familie. \cite{noe_burgen_online_ottenstein}

  \item[\textbf{Söhne und Enkel Hugos} (\textasciitilde{}1193--\textasciitilde{}1255)] Die Generation zwischen Hugo (zuletzt 1192) und den nächsten namentlich belegten Vertretern ist urkundlich nicht greifbar. Mehrere Urkunden nennen \enquote{die Herren von Ottenstein} als Zeugen, ohne Einzelpersonen zu identifizieren. Eine Linie blieb als Haupterben auf der Burg; eine weitere Nebenlinie führte das Prädiktat in abgekürzter Form weiter. Diese Rekonstruktion wurde aus Schenkungsurkunden an das Stift Zwettl erschlossen.

  \begin{description}
    \item[\textbf{Hadmar von Steinære} (\textasciitilde{}1180--\textasciitilde{}1238)] Jüngerer Sohn Hugos de Ottenstaine (Nebenlinienträger, nicht Erbe der Burg). Verwendete das verkürzte Prädiktat \textit{von Steinære} anstelle von \textit{von Ottenstaine}, wie es für nachgeborene Söhne üblich war. Verheiratet mit einer nicht namentlich bekannten Frau unbekannter Herkunft; Vater Kuonrâts von Steinære. (Fiktional -- Name Hadmar aus dem Ottensteiner Namensbestand; Konstrukt der Nebenlinie historisch plausibel)
    \begin{description}
      \item[\textbf{\gls{Kuonrat von Steinare}} (geb. \textasciitilde{}1213)] Nachgeborener Sohn, kein Erbanspruch auf Burg Ottenstein. Protagonist. (Fiktionale Figur)
    \end{description}
  \end{description}

  \item[\textbf{Konrad und Hadmar von Ottenstein} (\textasciitilde{}1260/70)] Zwei Brüder, die in Urkunden der zweiten Hälfte des 13.~Jahrhunderts gemeinsam genannt werden. Beide trugen den Beinamen \textit{Asinus} (lat. \enquote{Esel}) -- ein freiwillig gewählter Beiname, der im 13.~Jh. Demut symbolisieren sollte. Hadmar hatte die Erbtochter \textbf{Jutta}, Tochter des letzten Ottensteiner-Zweigs aus der Linie Rauheneck, geheiratet und nannte sich ebenfalls nach Ottenstein. Wahrscheinlich Söhne oder Enkel eines Bruders von Hadmar von Steinære (Hauptlinie). \cite{burgen_austria_ottenstein,schneider_ottenstein}

  \item[\textbf{Albero der Ottensteiner} (bezeugt 1346)] Wohnte noch auf der Burg. Wurde einige Jahre später Küchenmeister Herzogs Rudolfs~IV. von Österreich (\textit{Rudolf des Stifters}). \cite{burgen_austria_ottenstein}

  \item[\textbf{Erlöschen der Familie} (kurz nach 1411)] Kurz nach 1411 dürfte die Familie in männlicher Linie erloschen sein; die Burg gelangte spätestens 1446 an Tobias von Rohr. \cite{noe_burgen_online_ottenstein}
\end{description}

\medskip
\noindent\textit{Gesamtdauer:} Die Ottensteiner hielten die Burg als freies Eigen über fast drei Jahrhunderte (ca.~1160--nach 1411) -- eine für ein Ministerialengeschlecht ungewöhnlich lange, kontinuierliche Besitzfolge.

\subsubsection*{Retz, das Thayatal und die Grafschaft Hardegg}

An der Nordgrenze Niederösterreichs, wo die \textit{Thaya} eine tiefe Schlucht durch den Granit- und Gneis-Untergrund des Böhmischen Massivs gräbt, liegen drei historisch eng verbundene Orte: \textbf{Retz}, das \textbf{Thayatal} und \textbf{Hardegg}. \cite{kacetl_thayatal}

\textit{Retz} wurde 1180 erstmals als \textit{Rezze} urkundlich erwähnt -- ein slawisch geprägter Name, der vermutlich den nahen Retzbach bezeichnet. Die eigentliche \textit{Planstadt} Retz entstand erst im letzten Viertel des 13.~Jahrhunderts: Graf \textbf{Berthold von Rabenswalde}, dem \textbf{Rudolf~I. von Habsburg} nach 1278 die Grafschaft und Herrschaft \textbf{Hardegg} als Lehen verlieh, verlegte seinen Wohnsitz nach Retz. Er stiftete dort ein \textit{Dominikanerkloster} (Kirche 1295 fertiggestellt) und gründete um 1300 die reguläre Stadtanlage auf einem rechteckigen Grundriss in erhöhter Lage. \cite{kacetl_thayatal}

Das \textbf{Thayatal} -- die tiefe Durchbruchsschlucht der Thaya -- bildete die natürliche Grenze zwischen dem Herzogtum Österreich und dem böhmischen Mähren. Das steile Gelände mit seinen Felsvorsprüngen begünstigte die Anlage von Burgen, die die wichtigen Übergänge sicherten. Im Tal liegt die \textbf{Burg Hardegg}, Namensgeber der gleichnamigen Grafschaft und Burg der Grafen von Hardegg. Hardegg selbst ist eine der kleinsten Städte Österreichs und verdankt sein Stadtrecht dem mittelalterlichen Stellenwert des Übergangs über die Thaya. \cite{kacetl_thayatal,stenzel_burg_zu_burg}

Die \textit{Grafschaft Hardegg} gelangte 1314 aus dem Besitz der Käfernburger an das Geschlecht der \textbf{Querfurter} -- Linie Maidburg-Hardegg --, die sie bis zu ihrem Erlöschen 1483 innehatten. Im 13.~Jahrhundert war die gesamte Region zwischen Retz, Hardegg und dem Thayatal ein Nahtkontakt zwischen österreichischer und böhmisch-mährischer Herrschaft -- geprägt von der politischen Instabilität nach dem Aussterben der Babenberger und dem Aufstieg Ottokars~II. \cite{kacetl_thayatal,niederstaetter_oesterreich}

\subsubsection*{Die Sage vom Einsiedler im Thayatal}

Die \textbf{Einsiedlerhöhle} (auch \textit{Einsiedlerklause} genannt) klebt wie ein Schwalbennest etwa 5 bis 12~Meter hoch in einer fast senkrechten Felswand des Thayatals nahe der Stadt \textbf{Hardegg}. Um diese in den Fels gemauerte Behausung rankt sich eine tragische Sage aus der Zeit der Kreuzzüge. \cite{sagen_thayatal,kacetl_thayatal}

\paragraph*{Der grausamen Irrtum}

Ein Ritter verließ seine Heimat, um im Heiligen Land gegen die Ungläubigen zu kämpfen, und ließ seine junge Ehefrau und seinen erstgeborenen, noch kleinen Sohn zurück. Nach vielen Jahren voller Not und Gefahren ritt er seinen Truppen voraus, um so schnell wie möglich auf der heimischen Burg zu sein. Doch als er im Burghof ankam, sah er seine Frau, die einen jungen Mann zärtlich in den Armen hielt. Von rasender Eifersucht und blindem Zorn übermannt, zog er sein Schwert und stieß die Klinge mit solcher Wucht durch den Körper des Jünglings, dass sie auch seine Frau durchbohrte. Sterbend sank sie zu Boden und hauchte mit ihrem letzten Atemzug die furchtbare Wahrheit: \enquote{Jetzt hast du dein eigenes Kind ermordet!} Der junge Mann war kein Liebhaber gewesen, sondern sein eigener Sohn, der in den Jahren der Abwesenheit zu einem stattlichen jungen Mann herangewachsen war.

\paragraph*{Die Buße im Thayatal}

Von unerträglicher Reue und Verzweiflung erfüllt, verschenkte der Ritter seine Burg und all sein Hab und Gut. Er hüllte sich in ein kratziges Büßergewand und zog als ruheloser Pilger durch das Land. Nach langem Umherirren gelangte er schließlich in das damals noch wilde und völlig unwegsame Tal der Thaya nahe Hardegg. An einer Flussschleife versperrte ihm eine steil aufragende Felswand den Weg; in dieser Wand entdeckte er in luftiger Höhe eine kleine Felsnische. Er deutete dies als göttlichen Wink, kletterte hinauf und mauerte die Nische zu einer kleinen Klause um. Dort verbrachte er den Rest seines Lebens in völliger Einsamkeit als Einsiedler und wurde mit der Zeit als frommer und gütiger Mann in der ganzen Region bekannt.

Die gemauerte Felsbehausung existiert tatsächlich: Forschungen haben Spuren einer Feuerstelle und eines Rauchabzugs ergeben -- es haben dort in der Vergangenheit Menschen, vermutlich Eremiten, gehaust. Die Höhle ist heute ein beliebtes Wanderziel im \textbf{Nationalpark Thayatal} und über den blau markierten \textit{Einsiedlerweg} vom Nationalparkhaus in Hardegg erreichbar. Unterhalb der Höhle liegt die sogenannte \textit{Einsiedlerwiese} direkt an der Thaya. \cite{kacetl_thayatal}

\subsubsection*{Raabs an der Thaya und die Burg Raabs}\label{raabs}

An der Stelle, wo die \textbf{Deutsche Thaya} und die \textbf{Mährische Thaya} zusammenfließen, erhebt sich auf einem langgestreckten Felssporn die \textbf{Burg Raabs} -- eine der bedeutendsten Burganlagen des nördlichen Waldviertels und ein Ort, dessen Name bis heute das Bild Österreichs in der tschechischen Sprache prägt.

Die älteste schriftliche Quelle ist die \textit{Chronica Boemorum} des \textbf{Cosmas von Prag}: Er berichtet zum Oktober~1100, wie Herzog Bratislaus von Böhmen die Burg sechs Wochen lang belagerte, um den Přemysliden Lutold zur Herausgabe zu zwingen (\textit{castrum Racouz}). Der Name geht auf den althochdeutschen Personennamen \textit{Rātgōz} zurück, der von den slawischen Bewohnern Mährens und Böhmens zu \textit{Racouz} umgeformt wurde. Archäologische Befunde legen nahe, dass auf dem Burgberg bereits um 1050 ein früher Steinbau stand. \cite{hopf_raabs,noe_burgen_online_raabs,wiki_raabs}

\subsubsection*{\textit{Rakousko} -- wie Raabs zum Namen Österreichs wurde}

Die mährisch-böhmischen Nachbarn nannten die rund 50~km lange Grafschaft Raabs (von Raabs bis Litschau) im Mittelalter \textit{Rakoza} -- das \enquote{Raabser Land}. Diese Bezeichnung wurde sukzessive auf das Land \textit{hinter} Raabs, also das Herzogtum Österreich insgesamt, ausgedehnt. Heute bezeichnen die Tschechen ganz Österreich als \textbf{Rakousko}, und die Österreicher als \textbf{Rakušané} -- beides direkte Ableitungen des Burgnamens Raabs. \cite[S.~54]{stenzel_burg_zu_burg}\cite{wiki_raabs,kacetl_thayatal}

\subsubsection*{Die Grafen von Raabs und ihr Erbe}

\textbf{Gottfried~II. von Raabs} und sein Bruder \textbf{Konrad~I.} wurden 1105 von Kaiser Heinrich~IV. als Verantwortliche für die \textit{Nürnberger Burg} eingesetzt -- sie galten damit als die \textbf{ersten Burggrafen von Nürnberg}. Ihr Nachfahre \textbf{Gottfried~III.} ist der erste, bei dem der Titel \textit{burggravius de Norinberg} urkundlich nachweisbar ist. Die Raabser Grafen standen damit am Ursprung der Burggrafschaft Nürnberg -- und damit, über Generationen und Heiraten, am Ursprung des Hauses \textbf{Hohenzollern}: \cite{hopf_raabs,wiki_raabs}

Nach dem Erlöschen der Raabser Grafen im Mannesstamm (um~1191) teilte sich die Grafschaft durch die Heiraten der beiden Erbtöchter:\begin{itemize}
  \item Die älteste Tochter \textbf{Sophie} heiratete \textbf{Friedrich~III. von Zollern}, den Nürnberger Burggrafen -- sie wurde zur \enquote{Stammmutter} der \textbf{Hohenzollern} und damit der späteren preußischen Könige und deutschen Kaiser.
  \item Die Tochter \textbf{Agnes} heiratete Gebhard~III. von Dollnstein-Hirschberg; der westliche Teil mit Burg Raabs, Heidenreichstein und Litschau ging an das Geschlecht Hirschberg.
\end{itemize}

Um~1200 gelangte der Markt Raabs an Herzog \textbf{Leopold~VI. den Glorreichen}. Ende~1251 übergab \textbf{Ottokar~II.} die Herrschaft an die \textbf{Grafen von Plain-Hardegg}. Im Juni~1260 war die Burg im Besitz des böhmischen Adligen \textbf{Wok von Rosenberg}; sein Sohn \textbf{Heinrich von Rosenberg} musste Burg Raabs am 26.~März 1282 dem Habsburger Herzog \textbf{Albrecht~I.} abtreten. Ab~1297 war die gesamte Grafschaft landesfürstliches Lehen. \cite{wiki_burg_raabs,hopf_raabs}

\subsubsection*{Die Burg im 13. Jahrhundert}

Die Burg Raabs ist eine langgezogene Höhenburg, rund 200~m lang und maximal 40~m breit, durch das Gelände bedingt stark in die Länge gezogen. Die ältesten erhaltenen Bauteile -- Bergfried, romanische Burgkapelle (geweiht dem hl. Klemens, 12.~Jh.) und südlicher Bering -- stammen aus dem 12.~Jahrhundert. Bauuntersuchungen haben eine \enquote{Massivbebauung des 13.~Jahrhunderts} im Bereich oberhalb der \textit{Umkehr} (dem ersten Vorhof) nachgewiesen. Der Bergfried zeigt unten ein lagerhaftes Bruchsteinmauerwerk aus der Mitte des 12.~Jahrhunderts; ab ca.~10~m Höhe wechselt das Mauerwerk auf einen 7-eckigen Grundriss aus der zweiten Hälfte des 13.~Jahrhunderts -- also genau in der Zeitspanne, in der die Burg durch Rosenberger und später Habsburg-Hände ging. Im 13.~Jahrhundert war Raabs unmittelbares Grenzzeichen zwischen dem Herzogtum Österreich und dem böhmisch-mährischen Machtbereich. \cite{noe_burgen_online_raabs,wiki_burg_raabs}

\subsection{Ägypten und die Kreuzzüge des 13. Jahrhunderts}

Im 13.~Jahrhundert war \textit{Ägypten} das strategische Herzstück der muslimischen Welt. Die \textbf{Ayyubiden-Dynastie}, gegründet von \textbf{Saladin} (gest. 1193), herrschte von Kairo aus über Ägypten, Palästina, Syrien und große Teile der Levante. Die Kontrolle über Ägypten bedeutete Kontrolle über die Verbindungsrouten zwischen dem Mittelmeer und dem Roten Meer -- und damit auch über die Flanken des christlichen Königreichs Jerusalem. Daher richteten sich die großen Kreuzzüge des 13.~Jahrhunderts immer häufiger gegen Ägypten statt unmittelbar gegen das Heilige Land.

\subsubsection*{Der Kreuzzug von Damiette (1217--1221) und die österreichische Beteiligung}

Der \textbf{Fünfte Kreuzzug} (1217--1221) unternahm den ersten großen Angriff auf Ägypten. Unter der formellen Führung des päpstlichen Legaten \textbf{Pelagius von Albano} landeten die Kreuzfahrer 1218 an der Mündung des Nils und belagerten die Festungsstadt \textbf{Damiette}. Im November 1219 fiel die Stadt -- ein enormer Erfolg. Sultan \textbf{al-Kamil Muhammad} (reg. 1218--1238), der älteste Sohn des verstorbenen Sultans al-Adil, bot daraufhin den Kreuzfahrern Jerusalem im Tausch gegen Damiette an; Pelagius lehnte ab, da er auf die Ankunft Kaiser Friedrichs~II. wartete.

An diesem Kreuzzug nahm auch \textbf{Leopold~VI. der Glorreiche}, Babenberger Herzog von Österreich, teil. Die österreichischen Kontingente gelangten auf dem Seeweg nach Ägypten und kämpften in den Reihen des Kreuzfahrerheeres vor Damiette. Leopold~VI. war einer der aktivsten Kreuzzugsteilnehmer seiner Zeit -- er hatte bereits im Albigenserkreuzzug und in Spanien gekämpft. Sein Aufenthalt vor Damiette und in Ägypten dürfte erheblich zum Austausch zwischen österreichischer Ritterschaft und dem Orient beigetragen haben. \cite{lechner_babenberger}

Als die Kreuzfahrer 1221 ins Nildelta vordrangen, trieben sie sich in eine Falle: Das Nilhochwasser und die überlegene ägyptische Kavallerie schnitten den Rückzug ab. Die Kreuzfahrer kapitulierten; Damiette wurde den Ägyptern zurückgegeben. Der Kreuzzug war gescheitert, ohne bleibende Gebietsgewinne.

\subsubsection*{Der Kreuzzug Friedrichs II. und der Friede von Jaffa (1228--1229)}

Kaum einem anderen Kreuzzug haftet ein so eigentümlicher Charakter an wie dem des \textbf{Kaiser Friedrich~II.} Er brach 1228 ins Heilige Land auf -- exkommuniziert vom Papst, damit formal als Feind der Kirche. Trotzdem, oder vielleicht gerade deshalb, erzielte er das Ziel, das militärische Kreuzzüge nie erreicht hatten: die kampflose Rückgewinnung Jerusalems.

Friedrich~II. und Sultan \textbf{al-Kamil} kannten einander aus diplomatischem Briefwechsel seit Jahren. Beide schätzten Gelehrsamkeit, Wissenschaft und das Gespräch über Grenzen des Glaubens hinweg. Am \textbf{18.~Februar 1229} schlossen sie den \textbf{Frieden von Jaffa}: Jerusalem, Bethlehem, Nazaret und ein Verbindungskorridor zur Küste wurden für zehn Jahre an das Königreich Jerusalem abgetreten. Die Muslime behielten den Felsendom und die al-Aksa-Moschee und das Recht, diese zu besuchen; die Christen durften die Stadt nicht befestigen.

Der Vertrag wurde von beiden Seiten heftig kritisiert: Der Papst war empört, dass ein exkommunizierter Kaiser Erfolg hatte, wo fromme Kreuzfahrer gescheitert waren; muslimische Gelehrte und der Kalif in Bagdad verurteilten al-Kamil, weil er eine heilige Stadt preisgegeben hatte. Dennoch gilt der Friede von Jaffa als der einzige Kreuzzug, der sein primäres Ziel -- Jerusalem -- durch Diplomatie statt durch Blutvergießen gewann.

\subsubsection*{Der Sechste Kreuzzug: Ludwig IX. und der Ägyptenfeldzug (1248--1250)}

Jerusalem fiel 1244 erneut -- diesmal endgültig -- an muslimische Söldnertruppen (Chwarezmiya), die Sultan \textbf{as-Salih Ayyub} angeworben hatte. Dies löste den \textbf{Sechsten Kreuzzug} unter dem französischen König \textbf{Ludwig IX.} aus, der 1248 mit einem gewaltigen Heer aufbrach.

Die Strategie war dieselbe wie 1218: Erst Ägypten bezwingen, dann Jerusalem fordern. Am \textbf{4.~Juni 1249} erreichten die Kreuzfahrer die Nilmündung. Am folgenden Tag schlug Ludwig ein ägyptisches Heer und überquerte mit seinen Rittern das westliche Nilufer. Die Garnison von \textbf{Damiette} floh noch am selben Abend, und die Stadt fiel am \textbf{6.~Juni 1249} fast kampflos in christliche Hände -- ein Gegensatz zur jahrelangen Belagerung von 1218--1219.

Doch Ludwig wartete fünfeinhalb Monate in Damiette auf Verstärkungen, anstatt sofort nilaufwärts vorzustoßen. Sultan as-Salih, schwer erkrankt, stellte sein Hauptheer bei \textbf{al-Mansura} auf. Er starb am 22.~November 1249, ohne die Kreuzfahrer besiegt zu haben. Seine Regentin hielt den Tod geheim, bis der Thronfolger eintreffen konnte.

Am \textbf{8.~Februar 1250} überquerten die Kreuzfahrer den Nilarm bei al-Mansura. Ludwigs Bruder \textbf{Robert von Artois} führte die Vorhut -- und verfolgte entgegen dem ausdrücklichen Befehl des Königs die fliehenden Ägypter in die Stadt hinein. Dort gerieten die Ritter in einen Hinterhalt der \textbf{Mameluken} unter dem Emir \textbf{Rukn ad-Din Baibars}. In den engen Gassen konnten sich die schwere Kavallerie nicht formieren; fast die gesamte Vorhut, darunter Robert von Artois selbst und der Großmeister der Templer, wurde vernichtet.

Ludwig hielt das Hauptheer zusammen und begann die Belagerung al-Mansuras, doch die Position der Kreuzfahrer verschlechterte sich wochenlang: Krankheiten, Hunger, und die Blockade des Nils durch ägyptische Boote. Im \textbf{April 1250} begann der Rückzug. Baibars ließ das schwächelnde Heer einholen; bei \textbf{Fariskur} (6.~April 1250) wurden Ludwig~IX. und der Großteil seines Heeres gefangen genommen.

\subsubsection*{Vom Sultanat zum Mamlukensultan: der Aufstieg Baibars'}

Die unmittelbare Folge des Kreuzzuges war ein Machtwechsel in Ägypten. Sultan \textbf{Turan Schah}, der eben erst die Regentschaft übernommen hatte, wurde wenige Wochen nach der Gefangennahme Ludwigs von seinen eigenen Mameluken-Emiren ermordet (Mai 1250). Die Mameluken -- ursprünglich türkische Sklavensoldaten, die unter den Ayyubiden als Eliteeinheiten dienten -- übernahmen nun die Macht in Kairo.

Der entscheidende Mann war \textbf{Rukn ad-Din Baibars}, der Sieger von al-Mansura. Er stieg zum Sultan auf (reg. 1260--1277) und richtete die gesamte militärische Kraft des neuen Mamlukensultanats auf die systematische Zerstörung der verbleibenden Kreuzfahrerstaaten. Der Verlust Jerusalems 1244 und das Scheitern des Sechsten Kreuzzuges markierten das definitive Ende der Epoche, in der eine christliche Rückeroberung des Heiligen Landes militärisch realistisch erschien.

\subsection{Die Kaiser des Heiligen Römischen Reiches}

Das \textit{Heilige Römische Reich} war keine moderne Zentralmonarchie, sondern ein Geflecht aus Fürstentümern, Kirchenlehen und Reichsstädten, dessen Kaiser -- sofern er ein solcher war und nicht nur König -- von Papst und Fürsten zugleich abhing. Die Geschichte des Reiches im 13.~Jahrhundert ist geprägt von einem erbitterten Ringen zwischen Kaisertum und Papsttum sowie von der zunehmenden Verselbständigung der Territorialfürsten. \cite{niederstaetter_oesterreich}

\subsubsection*{Otto IV. (reg. 1209--1215)}

\textbf{Otto~IV.} aus dem Haus der \textit{Welfen} war der erste und einzige Welfen-Kaiser. Seine Krönung 1209 in Rom durch Papst Innozenz~III. galt zunächst als Triumph der päpstlichen Partei -- doch schon ein Jahr später brach der Konflikt aus: Otto versuchte, die Staufer-Güter in Unteritalien zurückzugewinnen und damit den Papst an seiner empfindlichsten Flanke zu treffen. Innozenz~III. bannte ihn 1210 und unterstützte den jungen \textbf{Friedrich~II.} als Gegenkönig. In der \textbf{Schlacht von Bouvines} (27.~Juli 1214) erlitt Otto~IV. eine vernichtende Niederlage gegen den französischen König Philipp~II. August -- das Ende seiner Herrschaft war damit besiegelt. Er starb 1218 ohne Nachkommen. \cite{niederstaetter_oesterreich}

\subsubsection*{Friedrich II. (reg. 1220--1250)}

Kein Herrscher des Mittelalters ist rätselhafter und faszinierender als \textbf{Friedrich~II.}, der \textit{Stupor Mundi} -- das Staunen der Welt. In Sizilien aufgewachsen, dreisprachig (Latein, Arabisch, Altfranzösisch), Förderer von Wissenschaft und Philosophie, sammelte er an seinem Hof in Palermo sowohl christliche als auch muslimische Gelehrte. Sein Werk \textit{De arte venandi cum avibus} (Über die Kunst des Falkenierens) gilt bis heute als frühe empirische Naturkunde.

Politisch war er zugleich König von Sizilien, König von Deutschland, König von Jerusalem -- und ab 1220 Kaiser des Heiligen Römischen Reiches. Sein Verhältnis zur Kirche war von chronischem Konflikt geprägt: Dreimal wurde er exkommuniziert. Papst Gregor~IX. nannte ihn öffentlich den \textit{Antichrist}; Friedrich antwortete mit zynischer Gelassenheit.

Für das Deutsche Reich bedeutete seine Herrschaft eine schleichende Aushöhlung der kaiserlichen Autorität: Mit der \textbf{Goldenen Bulle von Eger} (1213) und vor allem dem \textbf{Confoederatio cum principibus ecclesiasticis} (1220) und dem \textbf{Statutum in favorem principum} (1232) übertrug er den geistlichen und weltlichen Fürsten weitreichende Souveränitätsrechte -- im Tausch gegen ihre Unterstützung in seinen Kämpfen. Die Grundlage für den späteren deutschen Flickenteppich aus Einzelterritorien wurde in seiner Regierungszeit gelegt.

Friedrich~II. starb am \textbf{13.~Dezember 1250} in Apulien. Mit ihm endete die letzte Phase staufischer Kaisermacht. \cite{niederstaetter_oesterreich}

\subsubsection*{Das Interregnum (1250--1273)}\label{das_interregnum}

Mit dem Tod Friedrichs~II. begann das \textit{Interregnum} -- eine gut zwanzigjährige Periode ohne anerkannten deutschen König, die in der Geschichtsschreibung auch \textit{das Kaiserlose, die schreckliche Zeit} genannt wird. Die Fürsten hatten kein Interesse an einer starken Zentralgewalt und wählten entweder schwache Kandidaten oder ließen das Amt faktisch unbesetzt.

Friedrichs Sohn \textbf{Konrad~IV.} (reg. 1250--1254) konnte sich im Reich kaum behaupten; er starb 1254 in Italien an Malaria, gerade dreißigjährig. Sein Sohn \textbf{Konradin} -- das letzte legitime Glied der Staufer -- wurde 1268 nach der \textbf{Schlacht von Tagliacozzo} von Karl von Anjou gefangen genommen und am 29.~Oktober 1268 in Neapel enthauptet. Das Staufergeschlecht erlosch.

In Deutschland wurden unterdessen als \textit{Gegenkönige} gewählt: der thüringische Landgraf \textbf{Heinrich Raspe} (1246--1247, Papstkandidat), danach Graf \textbf{Wilhelm von Holland} (1247--1256) sowie ab 1257 gleich zwei Kandidaten ohne effektive Macht: der englische Graf \textbf{Richard von Cornwall} und der kastilische König \textbf{Alfons~X.} (\textit{der Weise}). Beide wurden von unterschiedlichen Kurfürstengruppen gewählt; keiner von beiden erschien je effektiv in Deutschland, keiner wurde je zum Kaiser gekrönt.

Während das Kaisertum vakant blieb, konsolidierten die Territorialfürsten ihre Macht. Straßenraub (das sogenannte \textit{Raubrittertum}) nahm zu, Städte schlossen eigene Bündnisse (\textit{Städtebünde}), und der Rechtsschutz für Reisende und Händler versiegte auf weiten Strecken. \cite{niederstaetter_oesterreich}

\subsubsection*{Rudolf I. von Habsburg (reg. 1273--1291)}

Die Wahl \textbf{Rudolfs~I. von Habsburg} am 1.~Oktober 1273 in Frankfurt beendete das Interregnum. Die Kurfürsten wählten ihn bewusst als \textit{mittleren} Kandidaten: mächtig genug, um Ordnung herzustellen, aber nicht so mächtig, dass er die Fürstensouveränität gefährdete. Der Habsburger war bis dahin ein Schweizer Grafenhaus ohne Königstitel -- kein Staufer, kein Welfe. \cite{niederstaetter_oesterreich}

Rudolfs größte Leistung war die Niederwerfung des böhmischen Königs \textbf{Ottokar~II. Přemysl}, der sich geweigert hatte, Österreich, Steiermark, Kärnten und Krain als Reichslehen anzuerkennen. In der \textbf{Schlacht auf dem Marchfeld} (26.~August 1278) besiegte Rudolf Ottokar, entscheidend unterstützt durch das Reiterheer König \textbf{Ladislaus'~IV. von Ungarn}; Ottokar fiel im Kampf. Rudolf belehnte seine Söhne mit Österreich und der Steiermark -- die \textit{Hausmacht Habsburg} im Osten war gegründet, und die Geschichte Österreichs hatte eine neue Richtung gefunden. \cite{niederstaetter_oesterreich}

Rudolf~I. wurde nie zum Kaiser gekrönt (der Papst verweigerte die Krönung), war aber als \textit{römisch-deutscher König} der effektivste Herrscher des Reiches seit Jahrzehnten. Nach seinem Tod 1291 versuchten die Kurfürsten wieder, die Habsburger zu übergehen -- erst 1438 wurde das Kaisertum dauerhaft habsburgisch. \cite{niederstaetter_oesterreich}

\section{Chronologische Übersicht}

Die folgende Liste versammelt alle im historischen Anhang genannten Ereignisse in zeitlicher Reihenfolge.

\begin{description}
  \item[\textbf{Okt.~1100}] Erste urkundliche Nennung der Burg Raabs (\textit{castrum Racouz}) in der Chronik des Cosmas von Prag.
  \item[\textbf{1105}] Gottfried~II. und Konrad~I. von Raabs als erste Burggrafen von Nürnberg eingesetzt.
  \item[\textbf{1177}] Erste urkundliche Erwähnung des \textit{Hugo de Ottenstaine} als Besitzer der Burg Ottenstein im Waldviertel.
  \item[\textbf{1180}] Erste urkundliche Erwähnung von Retz als \textit{Rezze}.
  \item[\textbf{1189--1191}] Gründung des Deutschen Ordens als Krankenpflegebruderschaft während der Belagerung von Akkon.
  \item[\textbf{\textasciitilde{}1191}] Erlöschen der Grafen von Raabs im Mannesstamm; Teilung der Grafschaft: Sophie von Raabs heiratet Friedrich~III. von Zollern (Stammmutter der Hohenzollern); westlicher Teil mit Burg Raabs an Hirschberg.
  \item[\textbf{19.~Feb.~1199}] Erhebung des Deutschen Ordens zum Ritterorden; Bestätigung durch Papst Innozenz~III.
  \item[\textbf{um~1200}] Markt Raabs gelangt an Herzog Leopold~VI. von Österreich.
  \item[\textbf{1202}] Gründung des Schwertbrüderordens in Riga.
  \item[\textbf{1209}] Kaiserkrönung Ottos~IV. in Rom.
  \item[\textbf{1210}] Exkommunikation Ottos~IV. durch Innozenz~III.; Friedrich~II. als Gegenkönig unterstützt.
  \item[\textbf{1213}] Goldene Bulle von Eger: Friedrich~II. überträgt den Reichsfürsten weitreichende Hoheitsrechte.
  \item[\textbf{27.~Jul.~1214}] Schlacht von Bouvines: vernichtende Niederlage Ottos~IV. gegen Philipp~II. von Frankreich; Ende seiner Herrschaft.
  \item[\textbf{1217}] Aufbruch zum Fünften Kreuzzug; Leopold~VI.~der Glorreiche nimmt mit österreichischen Kontingenten teil.
  \item[\textbf{1218}] Tod Ottos~IV.; Sultan al-Kamil übernimmt die Ayyubiden-Herrschaft; Kreuzfahrer beginnen Belagerung von Damiette.
  \item[\textbf{Nov.~1219}] Fall von Damiette.
  \item[\textbf{1220}] Kaiserkrönung Friedrichs~II.; \textit{Confoederatio cum principibus ecclesiasticis}.
  \item[\textbf{1221}] Kreuzfahrer im Nildelta eingekesselt; Kapitulation und Rückgabe Damiettes; Fünfter Kreuzzug gescheitert.
  \item[\textbf{1226}] Goldene Bulle von Rimini: Friedrich~II. sichert dem Deutschen Orden das Recht auf das künftige Prußenland.
  \item[\textbf{1228}] Friedrich~II. bricht exkommuniziert ins Heilige Land auf.
  \item[\textbf{18.~Feb.~1229}] Friede von Jaffa: Jerusalem, Bethlehem und Nazaret kampflos für zehn Jahre an das Königreich Jerusalem abgetreten.
  \item[\textbf{1230}] Tod Leopolds~VI.~des Glorreichen; Regierungsantritt Friedrichs~II.~des Streitbaren in Österreich; Vertrag von Kruschwitz (Kulmerland an den Deutschen Orden).
  \item[\textbf{1231}] Gründung der Ordensburg Thorn; Beginn der systematischen Prußenfeldzüge.
  \item[\textbf{1232}] \textit{Statutum in favorem principum}: weitere Entmachtung des Kaisertums zugunsten der weltlichen Fürsten.
  \item[\textbf{1236}] Niederlage des Schwertbrüderordens in der Schlacht von Schaulen; Verhängung der Reichsacht über Herzog Friedrich~II.~den Streitbaren.
  \item[\textbf{1237}] Union von Viterbo: Schwertbrüderorden wird als \textit{Livländischer Orden} in den Deutschen Orden eingegliedert.
  \item[\textbf{1238}] Tod Sultan al-Kamils.
  \item[\textbf{1240}] Ordensritter besetzen Pskow; Alexander Newski besiegt ein schwedisches Heer an der Newa.
  \item[\textbf{5.~Apr.~1242}] Niederlage des Deutschen Ordens auf dem Peipussee gegen Fürst Alexander Newski.
  \item[\textbf{Sommer~1242}] Friede mit Nowgorod; die Narwa wird als Ostgrenze des Ordensstaates festgelegt.
  \item[\textbf{1244}] Jerusalem fällt endgültig an chwaresmische Söldner Sultans as-Salih Ayyub; Ludwig~IX. legt Kreuzzuggelübde ab.
  \item[\textbf{15.~Jun.~1246}] Tod Friedrichs~II.~des Streitbaren in der Schlacht an der Leitha; Babenberger im Mannesstamm erloschen; Nachfolgekrise beginnt.
  \item[\textbf{1246--1247}] Landgraf Heinrich Raspe als päpstlicher Gegenkönig im Reich.
  \item[\textbf{1247}] Graf Wilhelm von Holland als Gegenkönig gewählt.
  \item[\textbf{1248}] Ludwig~IX. bricht mit einem gewaltigen Heer zum Sechsten Kreuzzug auf.
  \item[\textbf{4.--6.~Jun.~1249}] Kreuzfahrer landen in Ägypten; Damiette fällt fast kampflos.
  \item[\textbf{22.~Nov.~1249}] Tod Sultan as-Salihs; Geheimhaltung des Todes bis zur Ankunft des Thronfolgers.
  \item[\textbf{8.~Feb.~1250}] Überquerung des Nilarms bei al-Mansura; Hinterhalt der Mameluken unter Baibars; Vorhut vernichtet; Robert von Artois fällt.
  \item[\textbf{6.~Apr.~1250}] Niederlage bei Fariskur; Ludwig~IX. und der Großteil seines Heeres gefangen genommen.
  \item[\textbf{Mai~1250}] Ermordung Turan Schahs durch Mamelukenemire; Beginn des Mamlukensultanats.
  \item[\textbf{13.~Dez.~1250}] Tod Kaiser Friedrichs~II. in Apulien; Ende der staufischen Kaisermacht.
  \item[\textbf{1251}] Ottokar~II.~Přemysl marschiert in Österreich ein.
  \item[\textbf{1252}] Ottokar~II. heiratet Margarete von Österreich; wird formell Herzog von Österreich. Herrschaft Raabs gelangt an die Grafen von Plain-Hardegg.
  \item[\textbf{1253}] Ottokar~II. folgt seinem Vater auf dem böhmischen Thron.
  \item[\textbf{1254}] Tod Konrads~IV. in Italien; Interregnum vertieft sich.
  \item[\textbf{1257}] Doppelwahl: Richard von Cornwall und Alfons~X. von Kastilien als rivalisierende Gegenkönige.
  \item[\textbf{Jun.~1260}] Burg Raabs im Besitz des böhmischen Adligen Wok von Rosenberg.
  \item[\textbf{1260--1277}] Rukn ad-Din Baibars als Mamlukensultan; systematische Zerstörung der verbleibenden Kreuzfahrerstaaten.
  \item[\textbf{29.~Okt.~1268}] Hinrichtung Konradins in Neapel; Staufergeschlecht erloschen.
  \item[\textbf{1.~Okt.~1273}] Wahl Rudolfs~I. von Habsburg in Frankfurt; Ende des Interregnums.
  \item[\textbf{26.~Aug.~1278}] Schlacht auf dem Marchfeld: Niederlage und Tod Ottokars~II.; Österreich fällt an das Haus Habsburg.
  \item[\textbf{26.~März~1282}] Heinrich von Rosenberg übergibt Burg Raabs an den Habsburger Herzog Albrecht~I.
  \item[\textbf{1291}] Tod Rudolfs~I. von Habsburg.
  \item[\textbf{1295}] Fertigstellung der Dominikanerkirche in Retz (gestiftet von Graf Berthold von Rabenswalde).
  \item[\textbf{um~1300}] Stadtgründung Retz durch Graf Berthold von Rabenswalde.
\end{description}
